\subsection{The economic equilibrium}\label{\refsec:EE}

\smalltitle{With log preferences.}
Assume that we know the path of taxes $\tau_t$, pension characteristics $\theta_t, \epsilon_t$, and an initial exogenous level of savings $s_0$. In this case, recall that we simply have $B_{t+1}^i = \beta_{t,i}$. Given this, we can solve for the economic equilibrium without any numerical root finding: 
\begin{itemize}
    \item First compute $\Gamma_{s,t}$ from \eqref{\refeq:auxiliary:Gammas} (only defined for $t<T$).
    \item Then, compute finite horizon $\Theta_{h,t}$ from stacking \eqref{\refeq:auxiliary:Thetah} for $t<T$ and \eqref{\refeq:auxiliary:ThetahT} for $t=T$. 
    \item Then, compute $\Theta_{s,t}$ from \eqref{\refeq:auxiliary:Thetas} (only defined for $t<T$ as well).
    \item Given $s_0$ and $\Theta_{s,t}$ we can iterate forward on \eqref{\refeq:auxiliary:s} to solve for $s_t$ for all $t<T$. For $t=T$ we have $s_T = 0$ per construction.
\end{itemize}
As explained in \cref{\refmodelsec}, this set of core variables allows us to solve for individual labour supply and savings, factor prices, pension benefits, and finally consumption levels (from budgets).

We emphasize that the informal savings ratio $\iota_t$ is \emph{not} part of this core. Informal savings do not enter the formal capital stock, so none of $\Gamma_{s,t}, \Theta_{h,t}, \Theta_{s,t}, s_t, h_t$ depends on them, and the steps above are unaffected by the informal households' savings decision. Once the core is solved, $\iota_t$ follows in closed form from \eqref{\refeq:auxiliary:s0_s}, and is defined for $t<T$ only (at $t=T$ the informal, like the formal, save nothing). The same applies with CRRA preferences below, with $B_{t+1}^0$ evaluated at the solved $\left(s_t, h_{t+1}\right)$.




\smalltitle{With CRRA preferences.}
Assume that we know the path of taxes $\tau_t$, pension characteristics $\theta_t, \epsilon_t$, and an initial exogenous level of savings $s_0$. The core of the economic equilibrium can then be reduced to identifying vectors of $\Gamma_{s,t}, h_t, s_t$. The three vectors are identified by three sets of conditions (we have a square system of nonlinear equations):
\begin{itemize}
    \item The finite horizon labour supply condition requires \eqref{\refeq:auxiliary:h} with $\Theta_{h,t}$ being computed using \eqref{\refeq:auxiliary:Thetah} for $t<T$ and \eqref{\refeq:auxiliary:ThetahT} for $t=T$. 
    \item The finite horizon savings decision requires solving for $t<T$ (with $s_T = 0$); we use \eqref{\refeq:auxiliary:sFromH} for this. 
    \item Finally, $\Gamma_{s,t}$ is also only defined for $t<T$; we use \eqref{\refeq:auxiliary:Gammas} for this with condition \eqref{\refeq:auxiliary:Gammas} evaluated with $B_{t+1}^i(s_t, h_{t+1})$.
\end{itemize}
Given a solution to this core of the economic equilibrium ($s_t, h_t, \Gamma_{s,t}$), the rest of the outcomes can be solved as in the log case. 
    

\smalltitle{Initializing with steady state savings.}
When we solve for the economic equilibrium given vectors of $\tau_t, \theta_t$, and $\epsilon_t$, our default option for the initial exogenous level of savings $s_0$ is to solve for steady state levels using parameter values from the first period ($t=1$). 


With Log preferences, we have that $B^i = \beta_i$ and then we can simply solve for $\Gamma_{s}$ directly in closed-form using the general equation \eqref{\refeq:auxiliary:Gammas}. Given this, we can then compute steady state savings from \eqref{\refeq:steadystate_LOG:s} (here just repeated from \cref{\refmodelsec}):
\begin{subequations}\label{\refeq:steadystate_LOG}
    \begin{align}
        \Gamma_s =& \frac{1}{1+\xi}\frac{\sum_i \frac{\gamma_{i}\eta_{i}^{1+\xi}}{X_{i}^{\xi}}\frac{B^i}{1+B^i}}{1+\frac{1-\alpha}{\alpha}\frac{p\tau}{\kappa}\left(\theta+(1-\theta)\sum_i\frac{\gamma_{i}}{1+B^i}\right)} \label{\refeq:steadystate_LOG:Gammas}\\ 
         s^* =& \Gamma_h^{1+\xi}\left[\left(\frac{(1-\alpha)(1-\tau)}{\Gamma_h-\frac{1-\alpha}{\alpha}\frac{p\theta\tau}{\kappa}\Gamma_s}\right)^{1+\xi}\frac{\Gamma_s^{1+\alpha\xi}}{\nu^{\alpha(1+\xi)}}\right]^{\frac{1}{1-\alpha}} \label{\refeq:steadystate_LOG:s}
    \end{align}
\end{subequations}

With CRRA preferences, we have to identify $s^*$ numerically instead. This can be done technically in several ways, but we prefer the following as it allows us to solve the problem as a bounded scalar search where there are robust and fast algorithms available (e.g. golden-section search). Specifically, we solve for the $\Gamma_s$ that satisfies the following condition \eqref{\refeq:steadystate_CRRA:Gammas}
\begin{subequations}\label{\refeq:steadystate_CRRA}
\begin{align}
    &\Gamma_s(1+\xi)\left[1+\frac{1-\alpha}{\alpha}\frac{p\tau}{\kappa}\left(\theta+(1-\theta)\sum_i\frac{\gamma_i}{1+B^i(\Gamma_s)}\right)\right] = \sum_i\frac{\gamma_i\eta_i^{1+\xi}}{X_i^{\xi}} \frac{B^i(\Gamma_s)}{1+B^i(\Gamma_s)}, \label{\refeq:steadystate_CRRA:Gammas} \\
    &\text{where }B^i(\Gamma_s) = \left(\beta_i\right)^{\rho} \left(\frac{\alpha}{p}\frac{\Gamma_h-\frac{1-\alpha}{\alpha}\frac{p\theta\tau}{\kappa}\Gamma_s}{(1-\alpha)(1-\tau)}\frac{\nu}{\Gamma_s}   \right)^{\rho-1} \label{\refeq:steadystate_CRRA:Bi}
\end{align}
\end{subequations}
It is relatively straightforward to define a bounded interval of $\Gamma^s$ to test out. For instance, with constant $B^i$ across types, we have that $\Gamma_s \in \left(0, B/((1+B)(1+\xi)) \right)$, which means that we can practically always use a golden section search with bounds $(0, 0.9)$ and be sure to identify the solution.
